\section{Abbreviations and important terms}

\subsection{Abbreviations}

{\small
\begin{longtable}{@{}p{2.6cm}p{5.2cm}p{7.6cm}@{}}
\toprule
\textbf{Short form} & \textbf{Full form} & \textbf{Meaning in this project} \\
\midrule
\endhead
AoI & Age of Information & Time (in slots) since the sink last received a reading from a sensor. \\
AoCI & Age of Correlated Information & A sensor is only as stale as its freshest correlated neighbour; an early, graph-based version of our idea (kept as a metric). \\
AoII & Age of Incorrect Information & Variant of AoI that only counts time during which the sink's knowledge is wrong \cite{17}. \\
PAoI & Peak Age of Information & The largest age reached before an update arrives. \\
WSN & Wireless Sensor Network & Many small sensors reporting over radio to a base station. \\
IoT & Internet of Things & Networked everyday devices and sensors. \\
CPS & Cyber-Physical System & A system where computation monitors or controls a physical process. \\
GNN & Graph Neural Network & A neural network that works on graphs by passing messages between connected nodes. \\
GCN & Graph Convolutional Network & A common GNN type \cite{21}. \\
MPNN & Message-Passing Neural Network & General GNN framework; ours is one. \\
MLP & Multi-Layer Perceptron & A plain feed-forward neural network. \\
GELU & Gaussian Error Linear Unit & The activation function used in our MLPs. \\
LN & Layer Normalisation & Rescales each node's embedding to stabilise training. \\
RL & Reinforcement Learning & Learning a policy by trial and reward. \\
DRL & Deep Reinforcement Learning & RL where the policy is a neural network. \\
MDP & Markov Decision Process & The formal model of a sequential decision problem (state, action, reward, transition). \\
PPO & Proximal Policy Optimization & The RL algorithm we train with \cite{23}. \\
GAE & Generalized Advantage Estimation & How PPO estimates how much better an action was than expected. \\
KL & Kullback--Leibler divergence & Distance between two probability distributions; used to monitor PPO updates. \\
DQN & Deep Q-Network & A value-based RL algorithm (not used; see Section 7). \\
A2C & Advantage Actor-Critic & A simpler actor-critic RL algorithm (not used). \\
SAC & Soft Actor-Critic & An RL algorithm for continuous actions (not used). \\
BC & Behaviour Cloning & Training a network to imitate a teacher; tried and replaced. \\
KF & Kalman Filter & Optimal estimator for linear-Gaussian systems; our sink. \\
AR(1) & Autoregressive process of order 1 & Each reading depends on the previous one plus noise. \\
RBF & Radial Basis Function & The kernel that makes correlation fall off with distance in the simulator. \\
MSE & Mean Squared Error & Average squared estimation error over sensors and slots. \\
RMSE & Root Mean Squared Error & Square root of MSE, in physical units ($^\circ$C or mph). \\
CI & Confidence Interval & Range likely to contain the true mean (we use 95\%). \\
LS & Least Squares & Fitting method used for each sensor's $\rho$. \\
B1 & Baseline 1 & Our name for the Max-AoI scheduler. \\
v4 & Version 4 & The correlation-diverse heuristic from our mid-semester work. \\
KG & Kalman-greedy & Baseline that greedily reduces the sink's uncertainty. \\
PEMS & Performance Measurement System & Caltrans traffic data system; PEMS-BAY is its Bay-Area subset. \\
MAC & Medium Access Control & The network layer that decides who transmits when; scheduling lives here. \\
HARQ & Hybrid Automatic Repeat reQuest & Retransmission scheme on lossy links (future work). \\
D2D & Device-to-Device & Direct links between devices (mentioned in literature). \\
MWIS & Maximum Weighted Independent Set & NP-hard problem behind link scheduling (literature). \\
CPU / GPU & Central / Graphics Processing Unit & We train and run everything on a laptop CPU. \\
API & Application Programming Interface & The JSON endpoints the demo website calls. \\
HTTP & HyperText Transfer Protocol & How the browser talks to our local server. \\
JSON & JavaScript Object Notation & Text format for results and API messages. \\
CI (software) & Continuous Integration & GitHub Actions runs our tests on every push. \\
NPZ & NumPy zipped archive & File format of our processed datasets and fitted models. \\
\bottomrule
\end{longtable}
}

\subsection{Symbols}

{\small
\begin{longtable}{@{}p{2.6cm}p{12.8cm}@{}}
\toprule
\textbf{Symbol} & \textbf{Meaning} \\
\midrule
\endhead
$M$ & Number of sensors (45 for Intel Lab, 325 for PEMS-BAY, 50 in the main synthetic test). \\
$C$ & Channel capacity: the maximum number of sensors that may transmit in one slot. \\
$F$ & Freshness bound: no sensor's AoI may exceed $F$ slots. We use $F=\lfloor 3M/C\rfloor$. \\
$t$ & Time slot index (one slot = 5 minutes on real data). \\
$x_t$ & True readings of all $M$ sensors at slot $t$ (a vector). \\
$\hat x_t$ (or $m_t$) & The sink's estimate of $x_t$. \\
$y_t$ & The readings actually received in slot $t$ (only from polled sensors). \\
$S_t$ & The set of sensors polled in slot $t$. \\
$\mathcal F_t$ & Forced set: sensors that must be polled now to respect the freshness bound. \\
$a_i(t)$ & Age of Information of sensor $i$. \\
$\rho_i$ & Temporal correlation of sensor $i$: how much of its value carries over to the next slot. \\
$Q$ & Covariance matrix of the innovations (the new, unpredictable part of each step). \\
$P_t$ & The sink's posterior covariance: $P_{ii}$ is its uncertainty about sensor $i$, $P_{ij}$ the shared uncertainty. \\
$P_0$ & Stationary covariance: the sink's uncertainty when it knows nothing. \\
$r$ & Measurement noise variance ($10^{-3}$). \\
$g_i$ & Variance gain: how much total uncertainty disappears if sensor $i$ is polled next. \\
$\hat r_{ij}$ & Estimated correlation between the innovations of sensors $i$ and $j$. \\
$\Acorr$ & Correlation graph (edge weight $|\hat r_{ij}|$ where $|\hat r_{ij}|\ge 0.3$). \\
$\Apost$ & Posterior-correlation graph built from $P_t$ each slot. \\
$h_i$ & The GNN's embedding (64 numbers) for sensor $i$. \\
$d$ & Embedding size (64). \\
$f_\theta$ & The learned per-sensor correction in the decoder. \\
$G_{ij}$ & The learned pairwise interaction term in the decoder. \\
$\alpha$ & Weight on the Kalman gain term in the decoder (learnable, starts at 4). \\
$\pi_\theta$ & The policy (the actor network) with parameters $\theta$. \\
$V_\phi$ & The critic (value network) with parameters $\phi$. \\
$\gamma$ & Discount factor (0.95): how much future rewards count. \\
$\lambda$ & GAE parameter (0.9). \\
$\epsilon$ & PPO clip range (0.15). \\
$\hat A_t$ & Advantage estimate. \\
$\ell$ & Correlation length scale in the simulator. \\
\bottomrule
\end{longtable}
}

\subsection{Important and new terms}

{\small
\begin{longtable}{@{}p{3.6cm}p{11.8cm}@{}}
\toprule
\textbf{Term} & \textbf{Meaning} \\
\midrule
\endhead
Sink / base station & The central node that receives sensor readings and estimates the field. \\
Field & The physical quantity spread over space and time (temperature over a floor, speed over a road network). \\
Slot & One scheduling decision period. Real data: 5 minutes. \\
Polling & The sink asking a sensor to transmit in this slot. \\
Scheduler / policy & The rule that decides which sensors to poll. \\
Channel capacity & How many transmissions fit in one slot ($C$). \\
Interference / conflict & Two sensors too close to transmit in the same slot (only in synthetic tests). \\
Freshness guarantee & Hard rule: every sensor is polled at least every $F$ slots. \\
Freshness guard & The component that enforces it by forcing overdue sensors before the scheduler chooses. \\
Forced poll & A poll made because a sensor reached age $F$. \\
Starvation & A sensor that is never polled. Happens to greedy schedulers without the guard. \\
Spatial correlation & Nearby sensors read similar values. \\
Temporal correlation & A sensor's reading now is similar to its reading a slot ago. \\
Innovation & The part of a reading that the previous reading could not predict (the ``surprise''). \\
Correlation graph & Graph whose nodes are sensors and whose edges join sensors whose innovations move together. \\
Redundancy & Information the sink already has (or can infer) from other sensors. \\
Estimation error & Difference between the true reading and the sink's estimate. \\
Posterior & What the sink believes after using all received data. \\
Prior & What the sink believes before new data arrives. \\
Predict step & Kalman step that moves the estimate forward one slot and increases uncertainty. \\
Update step & Kalman step that uses received readings to correct the estimate and reduce uncertainty. \\
Kalman gain & How strongly a new reading changes each sensor's estimate. \\
Variance gain $g_i$ & Total uncertainty removed by polling sensor $i$. \\
Greedy & Choosing the best option for the current step only, without looking ahead. \\
Myopic & Same as greedy: optimising only the present slot. \\
Submodular & ``Diminishing returns'' property; makes greedy sensor selection near-optimal within one slot. \\
Oracle & A baseline that is told the true covariance; not deployable, used as a reference. \\
Ablation & Removing one component to measure what it contributes (here: message passing). \\
Embedding & A vector of numbers the network uses to describe one sensor. \\
Message passing & Each node updates its embedding using its neighbours' embeddings. \\
Permutation equivariance & Renumbering the sensors renumbers the output the same way; a GNN property. \\
Pointer decoder & A decoder that ``points at'' one input (sensor) at a time to build an ordered selection \cite{22}. \\
Autoregressive & Each pick depends on the picks already made. \\
Plackett--Luce & Probability model for an ordered list built by repeated softmax picks; gives exact log-probabilities. \\
Residual design & The network learns a correction on top of a known good rule (Kalman-greedy) instead of from scratch. \\
Softmax & Turns scores into probabilities that sum to 1. \\
Masking & Setting a score to $-\infty$ so an item cannot be picked. \\
Actor & The policy network (chooses actions). \\
Critic & The value network (predicts future reward; helps train the actor). \\
Reward & The number the agent tries to maximise: here, low estimation error. \\
Return & Discounted sum of future rewards. \\
Advantage & How much better an action did than the critic expected. \\
Rollout & A batch of experience collected with the current policy. \\
Entropy bonus & Small reward for keeping the policy random enough to explore. \\
Clipping & PPO's limit on how much one update can change the policy. \\
Warm start & Starting training from an earlier checkpoint. \\
Checkpoint & Saved network weights. \\
Validation set & Data used to choose the best checkpoint, never used for the final results. \\
Test set & Held-out data used only for final results. \\
Held-out / chronological split & Training uses earlier dates, testing uses later dates. \\
Data leakage & Accidentally letting test data influence training; prevented and tested. \\
Daily profile & The typical reading for each time of day, learned from training data. \\
Standardisation & Expressing values as deviations in units of standard deviation. \\
Recording gap / outage & A period where the dataset has no readings. \\
Burn-in & First slots of an episode that are not scored, so the sink can settle. \\
Episode & One run of a scheduler over a window of slots. \\
Test window & One disjoint stretch of the test period used as an episode. \\
Layout & One random placement of sensors (synthetic). \\
Uniform / clustered layout & Sensors spread evenly / grouped in clumps. \\
Length scale & How far correlation reaches in the simulator. \\
Pilot phase & Initial period where the sink hears everything, used to learn correlations (synthetic). \\
Win rate & Share of test windows or layouts where the GNN had lower error. \\
Wilcoxon signed-rank test & Paired statistical test of whether one method is consistently better. \\
p-value & Probability of seeing a difference this large by chance if the methods were equal. \\
Zero-shot & Tested on a setting never seen in training (e.g.\ 150 or 200 sensors). \\
Generalisation & Working well on unseen data or sizes. \\
Live demo & Our website that replays a real day and shows each scheduler's transmissions. \\
\bottomrule
\end{longtable}
}
